%% FullCourse_10Lecs -- Lecture 10, Topic 10.7: Synthesis + Appendix
%% Source: ./archive_pre_split/Lec10_Case_Studies_v2.tex (sections 7-8)

%% Shared preamble for the FullCourse_10Lecs topic decks.
%%
%% Each topic .tex \input's this file, then defines its own \title, \date
%% override (if any), \graphicspath, and \begin{document}.
%%
%% Reconstructed verbatim from the inlined copy preserved in
%% Lec10_Case_Studies/Lec10_T8_Case6_DeepRamsey.tex (lines 23-190), which is
%% explicitly annotated there as "from ../shared_preamble.tex, copied verbatim".
%% Base preamble matches MiniCourse_8hr/Slides/shared_preamble.tex; this version
%% additionally provides \sectiondivider, the conceptbox/cautionbox/examplebox/
%% takeawaybox tcolorboxes, and \compacttables used across the split decks.

\documentclass[10pt,english,aspectratio=169]{beamer}

\usepackage{ae,aecompl}
\usepackage[T1]{fontenc}
\usepackage{textcomp}
\usepackage[utf8]{inputenc}
\usepackage{babel}
\usepackage{datetime2}

\setcounter{secnumdepth}{3}
\setcounter{tocdepth}{3}

\usepackage{amsmath, amssymb, amsthm, graphicx}
\usepackage{booktabs, multirow}
\usepackage{tikz}
\usepackage{pgfplots}
\pgfplotsset{compat=1.16}
\usepackage[most]{tcolorbox}
\tcbuselibrary{skins, breakable, theorems}
\usepackage{listings}
\usepackage{xcolor}

\lstset{
  basicstyle=\ttfamily\small,
  keywordstyle=\color{blue!80!black}\bfseries,
  commentstyle=\color{green!50!black}\itshape,
  stringstyle=\color{orange!80!black},
  backgroundcolor=\color{gray!8},
  frame=single,
  rulecolor=\color{gray!40},
  breaklines=true,
  showstringspaces=false,
  numbers=none,
}

\lstdefinestyle{pythonstyle}{
    language=Python,
    basicstyle=\ttfamily\footnotesize,
    keywordstyle=\color{blue!80!black}\bfseries,
    commentstyle=\color{green!50!black}\itshape,
    stringstyle=\color{orange!80!black},
    frame=single,
    breaklines=true,
    showstringspaces=false,
    numbers=left,
    numbersep=5pt,
    numberstyle=\tiny\color{gray},
    backgroundcolor=\color{gray!8},
    rulecolor=\color{gray!40}
}

\ifx\hypersetup\undefined
  \AtBeginDocument{%
    \hypersetup{bookmarks=false,breaklinks=true,pdfborder={0 0 0},colorlinks=true,
      linkcolor=DarkText,urlcolor=RedTitle}%
  }
\else
  \hypersetup{bookmarks=false,breaklinks=true,pdfborder={0 0 0},colorlinks=true,
    linkcolor=DarkText,urlcolor=RedTitle}%
\fi

\makeatletter
\newcommand\makebeamertitle{%
  \begin{frame}[plain]
    \vspace*{1.0cm}
    \centering
    {\usebeamerfont{title}\usebeamercolor[fg]{title}\inserttitle\par}
    \vspace{1.0cm}
    {\usebeamerfont{author}\usebeamercolor[fg]{author}\insertauthor\par}
    \vspace{0.2cm}
    {\usebeamerfont{date}\usebeamercolor[fg]{date}\insertdate\par}
  \end{frame}}%
\AtBeginDocument{%
  \let\origtableofcontents=\tableofcontents
  \def\tableofcontents{\@ifnextchar[{\origtableofcontents}{\gobbletableofcontents}}
  \def\gobbletableofcontents#1{\origtableofcontents}%
}

\usetheme{metropolis}
\usefonttheme{professionalfonts}

\definecolor{LightGray}{RGB}{230,230,230}
\definecolor{RedTitle}{RGB}{0,0,200}
\definecolor{VeryLightGray}{RGB}{245,245,245}
\definecolor{DarkText}{RGB}{0,0,0}

\setbeamercolor{background canvas}{bg=white}
\setbeamercolor{title}{fg=RedTitle}
\setbeamercolor{author}{fg=DarkText}
\setbeamercolor{date}{fg=DarkText}
\setbeamercolor{frametitle}{bg=VeryLightGray,fg=DarkText}
\setbeamercolor{normal text}{fg=DarkText}
\setbeamerfont{title}{size=\large}

\setbeamersize{text margin left=5mm, text margin right=5mm}
\addtobeamertemplate{frametitle}{}{\vspace{-0.5em}}
\newcommand{\reducevspace}{\vspace{-0.5cm}}

% Shared section divider and box styles for split topic decks.
\newcommand{\sectiondivider}[2]{%
\begin{frame}[plain,noframenumbering]
\begin{center}
\vfill
{\Large\textbf{#1}\par}
\vspace{0.35cm}
{\normalsize #2\par}
\vfill
\end{center}
\end{frame}
}

\newtcolorbox{conceptbox}[1][]{
  colback=blue!3,
  colframe=RedTitle!55,
  boxrule=0.35mm,
  arc=1mm,
  left=2mm,
  right=2mm,
  top=1mm,
  bottom=1mm,
  #1
}
\newtcolorbox{cautionbox}[1][]{
  colback=red!3,
  colframe=red!50!black,
  boxrule=0.35mm,
  arc=1mm,
  left=2mm,
  right=2mm,
  top=1mm,
  bottom=1mm,
  #1
}
\newtcolorbox{examplebox}[1][]{
  colback=green!3,
  colframe=green!45!black,
  boxrule=0.35mm,
  arc=1mm,
  left=2mm,
  right=2mm,
  top=1mm,
  bottom=1mm,
  #1
}
\newtcolorbox{takeawaybox}[1][]{
  colback=VeryLightGray,
  colframe=RedTitle!55,
  boxrule=0.35mm,
  arc=1mm,
  left=2mm,
  right=2mm,
  top=1mm,
  bottom=1mm,
  #1
}
\newcommand{\compacttables}{\renewcommand{\arraystretch}{1.15}}

\setbeamertemplate{navigation symbols}{}
\setbeamertemplate{footline}{%
  \hfill%
  \usebeamercolor[fg]{page number in head/foot}%
  \usebeamerfont{page number in head/foot}%
  \insertframenumber\,/\,\inserttotalframenumber\kern1em\vskip2pt%
}

\makeatother

\author{Zhigang Feng}
% "date with time": compile timestamp so each build is identifiable.
\date{\today\ \DTMcurrenttime}


% Suppress redundant auto-generated metropolis section page; \sectiondivider frames are the dividers we keep.
\metroset{sectionpage=none}

\graphicspath{{./pic/}{./figures/}{./Exercise10_ES_Fellows/plots/}{../../MiniCourse_8hr/Slides/Module4_Agentic_AI/pic/}{../}}

\usetikzlibrary{positioning,arrows.meta,shapes,calc,fit,backgrounds}

\lstdefinestyle{markdownstyle}{
    basicstyle=\ttfamily\footnotesize,
    frame=single,
    rulecolor=\color{gray!40},
    backgroundcolor=\color{gray!8},
    breaklines=true,
    showstringspaces=false,
    numbers=none,
}

\title[AI for Econ Research]{AI for Economic Research: Dynamic Models, Language, and Agents\\
Lecture 10.7: Synthesis, Exercises, and Appendix}

\begin{document}

\makebeamertitle

%=============================================================================
\begin{frame}{Topic 10.7 Agenda}
\begin{enumerate}
  \item \textbf{Synthesis and Exercises}
  \medskip
  \item \textbf{Appendix}
\end{enumerate}
\end{frame}

%=============================================================================
\section{Synthesis and Exercises}
%===========================================================

\sectiondivider{Synthesis and Exercises}{Comparing reusable skills, empirical pipelines, public-web datasets, and full-lifecycle collaboration}

\begin{frame}{Compare the Case-Study Workflows}
\scriptsize
\begin{tabular}{p{2.3cm}p{3.0cm}p{3.2cm}p{3.1cm}}
\toprule
\textbf{Case} & \textbf{Artifact} & \textbf{Validation gate} & \textbf{Main failure mode} \\
\midrule
Paper review skill & referee-style report & quotes and section claims trace to the paper & generic criticism or hallucinated evidence \\
MEPS pipeline & panel, weights, figures & rates match external benchmarks & wrong crosswalk or dropped survey weights \\
ES Fellows & provenance-rich dataset & field-level source trace and missingness report & weak sources hidden by clean tables \\
Research collaborator & model, pipeline, critique, revision & claims survive adversarial review & polished output around weak framing \\
AI exposure measurement & rubric, JSON scores, comparison table, scatter & heuristic and reference comparison; invariant checks & construction choice silently drives coefficient sign \\
OLG homotopy sprint & V0--V5 version ladder and tests & each version inherits a validated baseline & one-shot model generation with no rollback point \\
\bottomrule
\end{tabular}
\end{frame}

\begin{frame}{When to Use Which Workflow}
\begin{columns}[T]
\begin{column}{0.24\textwidth}
\textbf{Build a skill}
\begin{itemize}\small
    \item repeated use
    \item high QC needs
    \item stable output format
\end{itemize}
\end{column}
\begin{column}{0.24\textwidth}
\textbf{Stay ad hoc}
\begin{itemize}\small
    \item one-off empirical job
    \item fast iteration
    \item small prompt set
\end{itemize}
\end{column}
\begin{column}{0.24\textwidth}
\textbf{Build a dataset harness}
\begin{itemize}\small
    \item many entities
    \item uneven sources
    \item provenance matters
\end{itemize}
\end{column}
\begin{column}{0.24\textwidth}
\textbf{Full-lifecycle collab.}
\begin{itemize}\small
    \item novel research question
    \item theory needs operationalization
    \item claims discipline required
    \item adversarial review needed
\end{itemize}
\end{column}
\end{columns}

\vspace{0.3cm}
\textbf{Measurement and computation cases.} Use the AI-exposure lab when the core risk is construct validity; use the OLG homotopy demo when the core risk is model complexity and rollback discipline.
\end{frame}

\begin{frame}{What Stayed Human Across All Cases}
\small
\begin{tabular}{p{5.9cm}p{6.0cm}}
\toprule
\textbf{Economist} & \textbf{Agent} \\
\midrule
define the question and success criteria & execute bounded tasks repeatedly \\
set the coding and validation rules & manipulate files, prompts, and reruns \\
spot implausible outputs & surface drafts, diagnostics, and plots \\
interpret the economic meaning & accelerate the path from question to artifact \\
\bottomrule
\end{tabular}

\vspace{0.3cm}
\textbf{The constant:} the question, identification logic, benchmark, and interpretation remain human work.
\end{frame}

\begin{frame}{Failure Modes to Watch For}
\begin{enumerate}
    \item \textbf{Case 1:} confident but unsupported criticism
    \item \textbf{Case 2:} clean-looking code with a wrong variable mapping
    \item \textbf{Case 3:} a rich-looking dataset built on weak or forbidden sources
    \item \textbf{Case 4:} AI validates the design instead of challenging it, or data contradicts theory without adversarial structure to surface it
    \item \textbf{All cases:} hidden assumptions that never made it into project files
\end{enumerate}

\vspace{0.3cm}
\textbf{The cure is the same in every case:} name the validation step, save the rules on disk, and inspect the output like a referee.
\end{frame}

\begin{frame}{Exercises and Deliverables}
\small
\begin{tabular}{p{2.6cm}p{4.3cm}p{5.7cm}}
\toprule
\textbf{Exercise} & \textbf{Artifact} & \textbf{Deliverable} \\
\midrule
Exercise A: MEPS lab & harmonized panel and validated figures & in-class guided run with the existing notebook \\
Exercise B: AI-exposure lab & rubric, scores, scatter, and validation table & run the Module 3 demo notebook; compare heuristic vs reference \\
Exercise C: OLG homotopy & V0->V1 transition or validation report & inspect the five-step demo; run one smoke test \\
Exercise D: fellows assignment & batch-level metadata dataset with provenance & normalized CSV, four figures, and a memo \\
\bottomrule
\end{tabular}

\vspace{0.3cm}
\begin{itemize}
    \item Exercise A teaches empirical execution and validation
    \item Exercise B teaches LLM-as-measurement and construct validation
    \item Exercise C teaches homotopy and rollback discipline for computational models
    \item Exercise D teaches controlled public-web data construction
\end{itemize}
\end{frame}



\begin{frame}{Capstone Demo Map}
\small
\begin{tabular}{p{2.6cm}p{4.2cm}p{5.0cm}}
\toprule
\textbf{Demo} & \textbf{Where to run} & \textbf{Expected visible output} \\
\midrule
AI exposure & Module 3 exposure demo notebook & substitution/assistance scatter; heuristic vs reference table; coefficient-flip discussion \\
OLG RAG & Module 4 OLG RAG terminal demo & retrieval report; top-$K$ evidence; grounded answer; refusal example \\
OLG homotopy & Module 4 OLG five-step notebook & V0--V5 version ladder; validation gates; regenerated reports \\
\bottomrule
\end{tabular}

\vspace{0.25cm}
{\footnotesize Full paths and commands are in \texttt{FullCourse\_Demo\_Index.md}.}

\vspace{0.15cm}
\textbf{Teaching rule:} every demo needs one launch command, one expected artifact, and one validation question students can answer from the output.
\end{frame}

\begin{frame}{Looking Back at the Course}
\begin{center}
\begin{tikzpicture}[
    node distance=0.15cm,
    block/.style={rectangle, rounded corners=3pt, draw=gray!50, fill=gray!5,
        minimum width=4.0cm, minimum height=0.55cm, font=\footnotesize, text=DarkText},
    current/.style={rectangle, rounded corners=3pt, draw=RedTitle, fill=blue!8,
        minimum width=4.0cm, minimum height=0.55cm, font=\footnotesize\bfseries, text=RedTitle},
]
\node[block] (l12) {Lec 1--6: models + AI methods};
\node[block, below=of l12] (l78) {Lec 7--8: LLMs + RAG};
\node[block, below=of l78] (l9) {Lec 9: agentic workflow};
\node[current, below=of l9] (l10) {Lec 10: case-study design};

\draw[-{Stealth[length=3mm]}, thick, RedTitle!40] (l12.south) -- (l78.north);
\draw[-{Stealth[length=3mm]}, thick, RedTitle!40] (l78.south) -- (l9.north);
\draw[-{Stealth[length=3mm]}, thick, RedTitle!40] (l9.south) -- (l10.north);
\end{tikzpicture}
\end{center}

\vspace{0.2cm}
\textbf{Final message:} the distance from question to artifact has collapsed, but economists still own the standards for what counts as evidence.
\end{frame}

%=============================================================================
\section{Appendix}
%===========================================================

\sectiondivider{Appendix}{Operational details moved out of the main narrative}

\begin{frame}{Appendix A: Paper-Review Skill Internals}
\small
\begin{tabular}{lp{9.2cm}}
\toprule
\textbf{Component} & \textbf{Role} \\
\midrule
\texttt{SKILL.md} & orchestrates the 7-phase pipeline \\
\texttt{shared-directives.md} & confidence gate, tone, quote rules \\
\texttt{section-prompts.md} & specialized prompts by section type \\
\texttt{editorial-prompt.md} & final quality gate before output \\
\bottomrule
\end{tabular}

\vspace{0.3cm}
\textbf{Reason this is appendix material:} students need to understand the workflow design before they need a file inventory.
\end{frame}

\begin{frame}[fragile]{Appendix B: MEPS Prompt Skeleton}
\begin{lstlisting}[style=pythonstyle, language={}, basicstyle=\ttfamily\scriptsize, numbers=none]
Prompt A: download + parse + harmonize + save panel
Prompt B: compute weighted rates by year and subgroup
Prompt C: generate annotated publication-quality figures
\end{lstlisting}

\vspace{0.25cm}
\textbf{Operational note:} the notebook remains the handout for the in-class lab.
\end{frame}

\begin{frame}{Appendix C: Fellows Normalized Schema}
\small
\begin{tabular}{p{3.6cm}p{8.1cm}}
\toprule
\textbf{Column group} & \textbf{Examples} \\
\midrule
seed fields & name, affiliation, election year \\
source fields & official profile URL, bio URL, CV URL \\
biography fields & birth date raw, birth year, PhD year, gender explicit \\
research fields & research fields raw, normalized field bucket \\
audit fields & source\_birth, source\_phd, source\_gender, source\_field, notes \\
\bottomrule
\end{tabular}
\end{frame}

\begin{frame}{Appendix D: Fellows Batching Rule}
\begin{itemize}
    \item official roster parsed into 18 batches of 50
    \item students complete one batch at a time
    \item every batch must report coverage before interpretation
    \item if birth-year coverage stays low, use years since PhD instead of age
\end{itemize}
\end{frame}

\begin{frame}{Appendix E: Resources}
\begin{itemize}
    \item Lecture 10 MEPS notebook in this folder
    \item \texttt{Exercise10\_ES\_Fellows/} script bundle in this folder
    \item Official Econometric Society current fellows page, checked April 13, 2026
    \item MEPS public-use files and external benchmark sources used in Case Study 2
\end{itemize}
\end{frame}

\end{document}
